\documentclass[12pt]{article}
\usepackage[T1]{fontenc}
\usepackage{lmodern}
\usepackage[margin=1in]{geometry}
\usepackage{booktabs}
\usepackage{setspace}
\setstretch{1.3}
\hyphenpenalty=10000
\exhyphenpenalty=10000
\sloppy
\pagestyle{plain}
\renewcommand{\thetable}{\arabic{table}}

\begin{document}

\thispagestyle{plain}
\begin{center}
{\Large\bfseries Learning to Launch: Alliance Experience and New Product Introductions in Manufacturing Firms\par}
\vspace{0.6em}
{\bfseries SYNTHETIC TEST PAPER, NOT REAL RESEARCH\par}
\vspace{0.6em}
Mara Quillon and Teodor Brask\\
Department of Imaginary Studies, Fictional University of Veltria
\end{center}

\begin{center}\bfseries Abstract\end{center}
\noindent How does a manufacturing firm's experience with strategic alliances shape its stream of new products? We study a panel of 412 manufacturing firms in Veltria, a fictional country, from 2008 to 2019. Our main finding is that alliance experience increases new product introductions (Hypothesis 1). Our second main finding is that organizational slack weakens the negative effect of product market competition on new product introductions (Hypothesis 7). Among our other hypothesized findings, supplier concentration has an inverted-U relationship with new product introductions (Hypothesis 12). We also find that the benefit of alliance experience grows with firm age, providing further support for our theory. The results show how collaborative learning builds the capability to launch.

\medskip
\noindent\textit{Keywords:} alliances; experience; new products; slack; competition

\noindent\textit{Note:} every firm, person, place, and number in this document is invented for testing a review tool.

\clearpage

\section{Introduction}

Manufacturing firms compete by renewing their product lines. A firm that brings new products to market more often can follow shifting demand, answer rivals, and spread the fixed costs of its plants over a wider range of goods. Yet the capability to launch new products is not evenly spread. Some firms renew their lines every year, while others with similar resources launch little. A long tradition in strategy explains such differences by what firms learn from their own history, and alliances are one of the main settings in which that learning takes place.

Alliances expose a firm to partners' technologies, routines, and markets. Each alliance also teaches the firm how to work with outsiders: how to pick a partner, divide tasks, share knowledge, and bring what was learned back into its own development teams. We argue that this accumulated alliance experience raises the rate at which a firm introduces new products, and we examine the conditions that shape that effect: the firm's resources, the competition it faces, its leadership and governance, and the structure of its supply base.

We test our arguments on a panel of 412 manufacturing firms in Veltria observed from 2008 to 2019. Veltria's national product register records every new product a manufacturer places on the domestic market, and its alliance register records every agreement that involves joint development, licensing, or joint marketing. The two registers let us follow each firm's alliance history and its stream of new products year by year.

The paper makes three contributions. First, it shows that alliance experience predicts new product introductions in a full panel of manufacturers rather than in a single industry. Second, it shows how slack and competition jointly shape the rate of introductions. Third, it examines how leadership, ownership, and governance condition the returns to alliance experience. We develop the hypotheses in Section 2, describe the data and measures in Section 3, report the results in Section 4, examine supplier concentration in Section 5, and discuss the findings in Section 6.

\section{Theory and Hypotheses}

\subsection{Alliance experience}

Experience with alliances builds a capability that is distinct from the knowledge any one partner provides. A firm that has formed many alliances has practiced the work of scanning for partners, writing agreements, setting up joint teams, and moving knowledge across the boundary of the firm. These routines lower the cost of the next collaboration and raise the share of partner knowledge that reaches the firm's own development teams. Alliance experience also widens the set of technical options a firm has seen, so that its engineers can recombine more ideas into new products. Both mechanisms point to a higher rate of new product introductions.

\textit{Hypothesis 1: Alliance experience is positively related to new product introductions.}

\subsection{Partner technological diversity}

The value of what a firm learns through alliances depends on whom it learns from. Partners whose technologies differ from one another expose the firm to a broader set of components, materials, and methods. When development teams draw on distant knowledge, they can combine ideas that rivals with a narrower set of partners never encounter.

\textit{Hypothesis 2: The technological diversity of a firm's alliance partners is positively related to new product introductions.}

\subsection{External CEO succession}

Leadership change can unlock new products. A chief executive appointed from outside the firm arrives without commitments to the existing product line and with ties to other firms and industries. Outside successors often revisit stalled development projects, reallocate budgets toward new lines, and bring in managers who favor launching products. Insiders, by contrast, tend to continue the strategies they helped build.

\textit{Hypothesis 3: External CEO succession increases new product introductions.}

\subsection{R\&D alliances}

Alliances differ in their purpose. Joint development and research agreements put the partners' engineers in the same teams and focus the collaboration on technologies that are close to launch. Licensing and marketing agreements move fewer of the routines that development requires. A firm whose alliance portfolio leans toward research and development therefore gains more of the knowledge that turns into products.

\textit{Hypothesis 4: The share of a firm's alliances that are R\&D alliances is positively related to new product introductions.}

\subsection{Product withdrawal}

Alliance experience should affect not only how many products a firm launches but also how well those products fare. Partners test designs, share market knowledge, and flag problems before launch. Firms with more alliance experience are better at absorbing such feedback, so their products should be less likely to be pulled from the market.

\textit{Hypothesis 5: Alliance experience lowers the risk that a new product is withdrawn from the market.}

\subsection{Process innovation}

The learning that alliances provide is specific to products: it concerns designs, features, and markets. Process innovation draws on a different body of knowledge, held mostly inside the plant and built through the firm's own production experience. We therefore expect alliance experience to have no bearing on process innovation, in either direction.

\textit{Hypothesis 6: Alliance experience is unrelated to process innovation.}

\subsection{Slack and competition}

Product market competition squeezes margins and shortens the planning horizon of managers. Under strong competition, firms cut development budgets and focus on defending existing lines, so competition tends to depress new product introductions. Organizational slack, in the form of resources beyond what current operations need, cushions the firm against this pressure. A firm with slack can keep development teams together and continue to fund launches when rivals press on prices. Slack thus weakens the negative effect of competition.

\textit{Hypothesis 7: Organizational slack weakens the negative effect of product market competition on new product introductions.}

\subsection{Founder CEOs}

Founders hold long-term views of their firms and have the standing to commit resources to collaborations that pay off slowly. A founder CEO can also overcome internal resistance to ideas that come from outside the firm. Both features make it easier for the firm to turn what it learned in alliances into new products.

\textit{Hypothesis 8: A founder CEO strengthens the positive relationship between alliance experience and new product introductions.}

\subsection{Family firms}

Family firms combine long horizons with dense personal ties to partners. These ties make it easier to sustain trust across many alliances and to absorb the tacit knowledge that partners share. We therefore expect family firms to gain more from each unit of alliance experience.

\textit{Hypothesis 9: The positive effect of alliance experience on new product introductions is stronger for family firms than for non-family firms.}

\subsection{Board independence}

Independent directors bring outside knowledge and monitoring. Outside directors sit on other boards, see other firms' product strategies, and press managers to keep the product line current. Boards with a larger share of independent directors should therefore push the firm toward more frequent launches.

\textit{Hypothesis 10: Board independence is positively related to new product introductions.}

\subsection{Absorptive capacity}

Alliance experience affects new products through the firm's ability to recognize, assimilate, and apply outside knowledge. Each alliance forces the firm to build the technical staff and routines needed to work with partners. That absorptive capacity, in turn, is what lets the firm convert outside knowledge into products.

\textit{Hypothesis 11: Absorptive capacity mediates the positive relationship between alliance experience and new product introductions.}

\begin{table}[h]
\centering\footnotesize
\textbf{Summary of Hypotheses}\par\smallskip
\begin{tabular}{lll}
\toprule
Hypothesis & Relationship & Predicted \\
\midrule
H1 & Alliance experience $\rightarrow$ new product introductions & Positive \\
H2 & Partner technological diversity $\rightarrow$ introductions & Positive \\
H3 & External CEO succession $\rightarrow$ introductions & Positive \\
H4 & R\&D alliance share $\rightarrow$ introductions & Positive \\
H5 & Alliance experience $\rightarrow$ risk of product withdrawal & Negative \\
H6 & Alliance experience $\rightarrow$ process innovation & No relationship \\
H7 & Slack moderates competition $\rightarrow$ introductions & Weakens negative effect \\
H8 & Founder CEO moderates alliance experience $\rightarrow$ introductions & Strengthens \\
H9 & Alliance experience $\rightarrow$ introductions, by family ownership & Stronger in family firms \\
H10 & Board independence $\rightarrow$ introductions & Positive \\
H11 & Alliance experience $\rightarrow$ absorptive capacity $\rightarrow$ introductions & Mediation \\
\bottomrule
\end{tabular}
\end{table}

\section{Methods}

\subsection{Data and sample}

Our sample consists of the 412 firms listed in the manufacturing index of the Veltrian business register at any time between 2008 and 2019. We follow each firm from 2008, or from its listing year, until 2019 or its delisting. The panel is unbalanced and holds 3,906 firm-years. Data on new products come from the Veltrian product register, which records each product a firm places on the domestic market and the date of any withdrawal. Data on alliances come from the national alliance register, which we read from 2001 so that every firm-year in the panel has a full five-year alliance history, also when alliance experience is lagged by two years. Financial data, ownership, board composition, and executive histories come from annual filings. The manufacturing index classifies firms by their registered activity, and some listed manufacturers draw most of their revenue from distribution or services; we keep them in the main sample and drop them in a robustness check.

\subsection{Dependent variables}

\textit{New product introductions.} Our main dependent variable is the number of new products the firm registered in the year. Because the count is skewed, the models use the log of one plus the count.

\textit{Industry rank in introductions.} As a second outcome, we rank the firms in each industry and year by their number of new products. Rank 1 is the firm with the most introductions, so a lower rank is better, and a variable that raises introductions relative to industry peers should have a negative coefficient.

\textit{Product withdrawal.} For Hypothesis 5, the unit of analysis is the product. We follow each of the 6,842 products introduced between 2008 and 2017 from its introduction until it is withdrawn or until the end of 2019. A total of 1,107 products were withdrawn.

\textit{Process innovation.} We measure process innovation as the log of one plus the number of new process filings the firm made with the national industrial property office in the year.

\textit{Absorptive capacity.} We measure absorptive capacity as the log of one plus the number of research and development staff the firm reports.

\subsection{Independent variables}

\textit{Alliance experience.} Alliance experience is the log of one plus the number of alliances the firm formed in the previous five years.

\textit{Partner technological diversity.} For each firm-year, we take the alliance partners of the previous five years and compute one minus the Herfindahl index of their primary technology classes.

\textit{R\&D alliance share.} The share of the firm's alliances in the previous five years that are joint development or joint research agreements.

\textit{Founder CEO and family firm.} Founder CEO equals one when the chief executive founded the firm. Family firm equals one when a family holds at least a quarter of the votes; no firm in the panel changes family status between 2008 and 2019.

\textit{External CEO succession.} External CEO succession equals one in the year a chief executive who had not worked at the firm before is appointed.

\textit{Board structure.} We measure board size as the number of directors on the board at the end of the fiscal year.

\textit{Supplier concentration.} Supplier concentration is the Herfindahl index of the firm's purchases across its suppliers. It ranges from 0.04 to 0.93; its 90th, 95th, 97th, and 99th percentiles are 0.55, 0.69, 0.78, and 0.86.

\textit{Slack and competition.} Slack is current assets minus current liabilities, scaled by total assets. Competition is one minus the Herfindahl index of sales in the firm's industry.

\subsection{Controls and estimation}

All firm-year models control for firm size (log employees), firm age (log years since founding), R\&D intensity (R\&D spending over sales), leverage, slack, and competition. Our main models are linear regressions of log introductions with firm and year fixed effects, and we call the model with all controls and both sets of fixed effects the full model. We estimate product withdrawal with Cox proportional hazards models stratified by industry. Interaction models include all constitutive terms (Brambor, Clark, \& Golder, 2006). Standard errors are clustered by firm. We use two-tailed tests throughout and treat p $<$ .05 as the criterion for support. Table 1 reports descriptive statistics.

\section{Results}

\subsection{Alliance experience}

Table 2 reports the tests of Hypothesis 1. Model 1 includes only the controls. Model 4 of Table 2, the full model, is our test of Hypothesis 1. In Model 4, alliance experience has a positive and significant coefficient (b = 0.142, SE = 0.041, p $<$ .01). Hypothesis 1 is supported. The estimate is smaller than in Models 2 and 3, which omit firm fixed effects, and close to the estimate in Model 5, which lags alliance experience by two years.

As a robustness check, Model 6 replaces the outcome with the firm's industry rank in introductions. Because rank 1 is the firm with the most new products, a lower rank is better, and a negative coefficient means that alliance experience moves the firm up the industry ranking. The coefficient is negative and significant (b = $-$0.315, SE = 0.098, p $<$ .01), in line with Hypothesis 1.

\subsection{Partner diversity}

Model 2 of Table 3 adds partner technological diversity and is our test of H2. Hypothesis 2 is supported (b = 0.061, SE = 0.052, p $<$ .05). The coefficient on alliance experience changes little when partner diversity enters the model.

\subsection{R\&D alliances}

Table 5 reports the tests of Hypothesis 4; column 1, the full model on the full sample, is the test of H4. Columns 2 to 4 re-estimate it on 2010 to 2019, without the largest firms, and on manufacturing-only firms. Across the four columns, the coefficient on R\&D alliance share is positive. H4 received partial support.

\subsection{Product withdrawal}

Hypothesis 5 predicts that alliance experience lowers the risk that a new product is withdrawn. Model 2 of Table 6 reports a Cox model of the withdrawal hazard. Alliance experience is significantly related to withdrawal (b = 0.18, SE = 0.06, p $<$ .01), supporting H5.

\subsection{Process innovation}

Hypothesis 6 states that alliance experience is unrelated to process innovation. In Model 2 of Table 7, the coefficient on alliance experience is small and not significant (b = 0.012, SE = 0.020). H6 is not supported.

\subsection{Slack and competition}

Hypothesis 7 predicts that slack weakens the negative effect of competition on new product introductions. Table 4 reports the tests. In Model 3, competition has a negative coefficient ($-$0.25, SE = 0.06, p $<$ .01), and the product term Slack $\times$ Competition is negative and significant (b = $-$0.090, SE = 0.030, p $<$ .01), supporting H7. The product term keeps its size when the sample is restricted to 2010 to 2019 (Model 4).

Model 5 of Table 4 adds the product of firm age and alliance experience. The coefficient is positive and significant (b = 0.058, SE = 0.021, p $<$ .01): older firms gain more from their alliance experience, providing further support for our theory.

\subsection{Founder CEOs}

Table 8 tests Hypothesis 8. Column 1 adds a founder CEO indicator to the full model, and column 2 adds its product with alliance experience. As Hypothesis 8 predicts, a founder CEO weakens the positive association between alliance experience and introductions: the product term in column 2 is negative and significant (b = $-$0.049, SE = 0.025, p $<$ .05). Hypothesis 8 is therefore supported.

\subsection{Family firms}

Hypothesis 9 predicts that the effect of alliance experience is stronger in family firms. Table 9 estimates the full model separately for the 136 family firms and the 276 non-family firms. For family firms, the coefficient on alliance experience is 0.18 (SE = 0.07, p $<$ .05); for non-family firms it is 0.09 (SE = 0.08) and not significant. Because the effect is significant for family firms but not for non-family firms, H9 is supported.

\subsection{Board independence}

Table 10 reports the tests of Hypothesis 10. Model 4 of Table 10, the full model, is our test of Hypothesis 10. Model 2 enters board size without controls, and Model 3 adds the controls without firm fixed effects. In Model 2, board size has a positive and significant coefficient (b = 0.083, SE = 0.031, p $<$ .01); Hypothesis 10 is supported.

\subsection{Absorptive capacity}

Table 11 tests Hypothesis 11 with pooled models that include industry and year fixed effects, because absorptive capacity changes slowly within firms. Following Baron and Kenny (1986), we proceed in three steps. In Model 1, alliance experience predicts introductions (b = 0.30, SE = 0.08, p $<$ .01). In Model 2, alliance experience predicts absorptive capacity (b = 0.42, SE = 0.09, p $<$ .01). In Model 3, absorptive capacity predicts introductions with alliance experience in the model (b = 0.51, SE = 0.12, p $<$ .01). The coefficient on alliance experience falls from 0.30 to 0.08 and is no longer significant, indicating full mediation (H11 supported).

\section{Supplier Concentration}

While estimating the models above, we examined how the structure of the supply base relates to new products. Firms that buy from a few suppliers can coordinate closely with them on new components, but firms whose purchases are concentrated in very few suppliers depend on those suppliers' own development plans and face hold-up when they want to change a design. Given the unexpected negative effect of supplier concentration at high levels, we propose H12:

\textit{Hypothesis 12: Supplier concentration has an inverted-U relationship with new product introductions.}

Table 12 reports the test. Model 2 adds supplier concentration and its square to the full model. The squared term is negative and significant (b = $-$1.25, SE = 0.38, p $<$ .01), supporting H12. The linear term is positive (b = 1.94, SE = 0.52). Moderate supplier concentration is associated with the most introductions.

\section{Discussion}

We set out to explain why some manufacturing firms launch new products more often than others. In a panel of Veltrian manufacturers, alliance experience predicts more new product introductions, and slack weakens the negative effect of competition on introductions. These two findings are the core of our contribution. They suggest that the capability to launch is learned through collaboration and protected by resources that let firms keep developing products when competition tightens.

Several of our other hypotheses concern who leads and owns the firm. Our evidence on external CEO succession (Hypothesis 3) is suggestive: outside successors are rare in the panel, appearing in 4 percent of firm-years, and the firms that appointed them were often in the middle of renewing their product lines. Leadership change thus seems to work alongside the learning that alliances provide.

\subsection{Post hoc exploratory analysis}

The analysis in this subsection was not hypothesized, and we report it as exploratory. After the main analyses were complete, we asked whether the location of partners matters. Adding to the full model the share of the firm's alliances formed with partners based outside Veltria, we find a positive coefficient (b = 0.047, SE = 0.019, p $<$ .05). Partners abroad may bring knowledge that domestic partners lack. Because we did not predict this result, it should be tested in new data before it is given weight.

\subsection{Limitations}

Our study has several limitations. Alliance experience is not randomly assigned, and firms that form many alliances may differ from others in ways our fixed effects do not absorb. Our measure of new products counts registrations and does not weigh products by their novelty or sales. The Veltrian registers cover one national market, so the results may not travel to economies with different alliance regulation. Future work could use variation in the timing of alliance formation to sharpen the causal interpretation.

\subsection{Conclusion}

Firms learn to launch. Experience with alliances builds routines that turn outside knowledge into new products, and slack helps firms keep launching when competition intensifies. Managers who treat alliances as a source of lasting capability, rather than as a series of one-off deals, can raise the pace at which their firms renew their product lines.

\section*{Appendix: Additional Robustness Checks}

We re-estimated the main models in three ways. First, measuring alliance experience over three years instead of five gives an estimate for Hypothesis 1 of 0.131 (SE = 0.040) in the full model. Second, dropping the 2008 and 2009 crisis years leaves the Slack $\times$ Competition product term almost unchanged ($-$0.087, SE = 0.032), as Model 4 of Table 4 shows. Third, estimating the partner diversity model with a Poisson count model of introductions gives a coefficient of 0.058 (SE = 0.049). These checks cover the alliance, slack, and partner diversity hypotheses.

\section*{References}

\begin{list}{}{\leftmargin=2em \itemindent=-2em \itemsep=2pt}
\item Baron, R. M., \& Kenny, D. A. 1986. The moderator-mediator variable distinction in social psychological research: Conceptual, strategic, and statistical considerations. \textit{Journal of Personality and Social Psychology}, 51(6): 1173--1182.
\item Brambor, T., Clark, W. R., \& Golder, M. 2006. Understanding interaction models: Improving empirical analyses. \textit{Political Analysis}, 14(1): 63--82.
\end{list}

\clearpage
\begin{table}[p]
\centering\small
\caption{Descriptive Statistics}
\begin{tabular}{lcccc}
\toprule
Variable & Mean & SD & Min & Max \\
\midrule
New product introductions (count) & 3.42 & 4.87 & 0 & 41 \\
Industry rank in introductions & 8.61 & 6.14 & 1 & 31 \\
Process innovation (log) & 0.61 & 0.74 & 0 & 3.91 \\
Alliance experience (log) & 1.38 & 0.97 & 0 & 4.62 \\
Partner technological diversity & 0.44 & 0.21 & 0 & 0.91 \\
R\&D alliance share & 0.37 & 0.29 & 0 & 1 \\
External CEO succession & 0.04 & 0.20 & 0 & 1 \\
Founder CEO & 0.18 & 0.38 & 0 & 1 \\
Family firm & 0.33 & 0.47 & 0 & 1 \\
Board size (directors) & 8.71 & 2.34 & 4 & 17 \\
Absorptive capacity (log R\&D staff) & 3.08 & 1.22 & 0 & 6.95 \\
Supplier concentration & 0.31 & 0.17 & 0.04 & 0.93 \\
Firm size (log employees) & 6.12 & 1.41 & 2.30 & 10.83 \\
Firm age (log years) & 3.21 & 0.62 & 1.10 & 4.80 \\
R\&D intensity & 0.034 & 0.029 & 0 & 0.214 \\
Leverage & 0.27 & 0.16 & 0 & 0.88 \\
Slack & 0.19 & 0.12 & 0.01 & 0.74 \\
Competition & 0.71 & 0.18 & 0.12 & 0.98 \\
\bottomrule
\end{tabular}
\begin{minipage}{0.92\textwidth}\footnotesize\smallskip
\textit{Note.} 3,906 firm-years from 412 Veltrian manufacturing firms, 2008 to 2019. Industry rank: 1 is the firm with the most introductions in its industry and year. Alliance experience is the log of one plus the number of alliances the firm formed in the previous five years.
\end{minipage}
\end{table}
\clearpage

\begin{table}[p]
\centering\small
\caption{Alliance Experience and New Product Introductions (Hypothesis 1)}
\begin{tabular}{lcccccc}
\toprule
 & Model 1 & Model 2 & Model 3 & Model 4 & Model 5 & Model 6 \\
 & Intro. & Intro. & Intro. & Intro. & Intro. & Rank \\
\midrule
Alliance experience &  & 0.231$^{***}$ & 0.188$^{***}$ & 0.142$^{***}$ & 0.121$^{***}$ & $-$0.315$^{***}$ \\
 &  & (0.038) & (0.040) & (0.041) & (0.043) & (0.098) \\[2pt]
Firm size & 0.118$^{***}$ &  & 0.104$^{***}$ & 0.097$^{***}$ & 0.101$^{***}$ & $-$0.274$^{***}$ \\
 & (0.031) &  & (0.030) & (0.032) & (0.032) & (0.081) \\[2pt]
Firm age & $-$0.052 &  & $-$0.061 & $-$0.048 & $-$0.050 & 0.112 \\
 & (0.044) &  & (0.039) & (0.045) & (0.045) & (0.121) \\[2pt]
R\&D intensity & 2.84$^{***}$ &  & 2.41$^{***}$ & 2.29$^{***}$ & 2.33$^{***}$ & $-$6.02$^{***}$ \\
 & (0.71) &  & (0.66) & (0.69) & (0.70) & (1.88) \\[2pt]
Leverage & $-$0.143 &  & $-$0.128 & $-$0.121 & $-$0.125 & 0.288 \\
 & (0.092) &  & (0.085) & (0.090) & (0.091) & (0.241) \\[2pt]
Slack & 0.231$^{**}$ &  & 0.196$^{**}$ & 0.188$^{*}$ & 0.192$^{**}$ & $-$0.402 \\
 & (0.098) &  & (0.094) & (0.097) & (0.097) & (0.261) \\[2pt]
Competition & $-$0.226$^{***}$ &  & $-$0.205$^{***}$ & $-$0.198$^{***}$ & $-$0.201$^{***}$ & 0.511$^{***}$ \\
 & (0.071) &  & (0.066) & (0.070) & (0.071) & (0.188) \\[2pt]
\midrule
Firm fixed effects & Yes & No & No & Yes & Yes & Yes \\
Year fixed effects & Yes & Yes & Yes & Yes & Yes & Yes \\
Observations & 3,906 & 3,906 & 3,906 & 3,906 & 3,906 & 3,906 \\
R$^2$ & 0.081 & 0.034 & 0.102 & 0.094 & 0.091 & 0.077 \\
\bottomrule
\end{tabular}
\begin{minipage}{0.92\textwidth}\footnotesize\smallskip
\textit{Note.} Dependent variable in Models 1 to 5: log of one plus new product introductions; in Model 6: industry rank in introductions, where 1 is the firm with the most introductions, so a lower rank is better. Model 4 is the full model. Model 5 lags alliance experience by two years. Standard errors, clustered by firm, in parentheses. \textasteriskcentered{} p$<$0.10, \textasteriskcentered{}\textasteriskcentered{} p$<$0.05, \textasteriskcentered{}\textasteriskcentered{}\textasteriskcentered{} p$<$0.01
\end{minipage}
\end{table}
\clearpage

\begin{table}[p]
\centering\small
\caption{Partner Technological Diversity and New Product Introductions (Hypothesis 2)}
\begin{tabular}{lcc}
\toprule
 & Model 1 & Model 2 \\
\midrule
Partner technological diversity &  & 0.061 \\
 &  & (0.052) \\[2pt]
Alliance experience & 0.142$^{***}$ & 0.138$^{***}$ \\
 & (0.041) & (0.041) \\[2pt]
Firm size & 0.097$^{***}$ & 0.096$^{***}$ \\
 & (0.032) & (0.032) \\[2pt]
Firm age & $-$0.048 & $-$0.047 \\
 & (0.045) & (0.045) \\[2pt]
R\&D intensity & 2.29$^{***}$ & 2.25$^{***}$ \\
 & (0.69) & (0.69) \\[2pt]
Leverage & $-$0.121 & $-$0.119 \\
 & (0.090) & (0.090) \\[2pt]
Slack & 0.188$^{*}$ & 0.186$^{*}$ \\
 & (0.097) & (0.097) \\[2pt]
Competition & $-$0.198$^{***}$ & $-$0.196$^{***}$ \\
 & (0.070) & (0.070) \\[2pt]
\midrule
Firm fixed effects & Yes & Yes \\
Year fixed effects & Yes & Yes \\
Observations & 3,906 & 3,906 \\
R$^2$ & 0.094 & 0.095 \\
\bottomrule
\end{tabular}
\begin{minipage}{0.92\textwidth}\footnotesize\smallskip
\textit{Note.} Dependent variable: log of one plus new product introductions. Standard errors, clustered by firm, in parentheses. \textasteriskcentered{} p$<$0.10, \textasteriskcentered{}\textasteriskcentered{} p$<$0.05, \textasteriskcentered{}\textasteriskcentered{}\textasteriskcentered{} p$<$0.01
\end{minipage}
\end{table}
\clearpage

\begin{table}[p]
\centering\small
\caption{Slack, Competition, and New Product Introductions (Hypothesis 7)}
\begin{tabular}{lccccc}
\toprule
 & Model 1 & Model 2 & Model 3 & Model 4 & Model 5 \\
\midrule
Alliance experience & 0.142$^{***}$ & 0.142$^{***}$ & 0.137$^{***}$ & 0.129$^{***}$ & 0.136$^{***}$ \\
 & (0.041) & (0.041) & (0.041) & (0.044) & (0.041) \\[2pt]
Slack &  & 0.188$^{*}$ & 0.204$^{**}$ & 0.211$^{**}$ & 0.201$^{**}$ \\
 &  & (0.097) & (0.096) & (0.101) & (0.096) \\[2pt]
Competition &  & $-$0.20$^{***}$ & $-$0.25$^{***}$ & $-$0.26$^{***}$ & $-$0.24$^{***}$ \\
 &  & (0.07) & (0.06) & (0.07) & (0.06) \\[2pt]
Slack $\times$ Competition &  &  & $-$0.090$^{***}$ & $-$0.087$^{***}$ & $-$0.088$^{***}$ \\
 &  &  & (0.030) & (0.032) & (0.030) \\[2pt]
Firm age $\times$ Alliance experience &  &  &  &  & 0.058$^{***}$ \\
 &  &  &  &  & (0.021) \\[2pt]
Firm size & 0.099$^{***}$ & 0.097$^{***}$ & 0.096$^{***}$ & 0.102$^{***}$ & 0.095$^{***}$ \\
 & (0.032) & (0.032) & (0.032) & (0.034) & (0.032) \\[2pt]
Firm age & $-$0.051 & $-$0.048 & $-$0.046 & $-$0.044 & $-$0.046 \\
 & (0.045) & (0.045) & (0.045) & (0.048) & (0.045) \\[2pt]
R\&D intensity & 2.36$^{***}$ & 2.29$^{***}$ & 2.27$^{***}$ & 2.19$^{***}$ & 2.24$^{***}$ \\
 & (0.69) & (0.69) & (0.69) & (0.73) & (0.69) \\[2pt]
Leverage & $-$0.131 & $-$0.121 & $-$0.118 & $-$0.109 & $-$0.117 \\
 & (0.090) & (0.090) & (0.090) & (0.094) & (0.090) \\[2pt]
\midrule
Firm fixed effects & Yes & Yes & Yes & Yes & Yes \\
Year fixed effects & Yes & Yes & Yes & Yes & Yes \\
Observations & 3,906 & 3,906 & 3,906 & 3,294 & 3,906 \\
R$^2$ & 0.088 & 0.094 & 0.099 & 0.101 & 0.103 \\
\bottomrule
\end{tabular}
\begin{minipage}{0.92\textwidth}\footnotesize\smallskip
\textit{Note.} Dependent variable: log of one plus new product introductions. Model 1 omits slack and competition. Model 4 re-estimates Model 3 on 2010 to 2019. Slack, competition, and firm age are mean-centered before the product terms are formed. Standard errors, clustered by firm, in parentheses. \textasteriskcentered{} p$<$0.10, \textasteriskcentered{}\textasteriskcentered{} p$<$0.05, \textasteriskcentered{}\textasteriskcentered{}\textasteriskcentered{} p$<$0.01
\end{minipage}
\end{table}
\clearpage

\begin{table}[p]
\centering\small
\caption{R\&D Alliance Share and New Product Introductions (Hypothesis 4)}
\begin{tabular}{lcccc}
\toprule
 & (1) & (2) & (3) & (4) \\
\midrule
R\&D alliance share & 0.034 & 0.039 & 0.029 & 0.071$^{**}$ \\
 & (0.029) & (0.031) & (0.033) & (0.030) \\[2pt]
Alliance experience & 0.140$^{***}$ & 0.133$^{***}$ & 0.146$^{***}$ & 0.151$^{***}$ \\
 & (0.041) & (0.043) & (0.045) & (0.044) \\[2pt]
Firm size & 0.096$^{***}$ & 0.101$^{***}$ & 0.094$^{***}$ & 0.093$^{***}$ \\
 & (0.032) & (0.034) & (0.034) & (0.034) \\[2pt]
Firm age & $-$0.047 & $-$0.044 & $-$0.050 & $-$0.039 \\
 & (0.045) & (0.048) & (0.049) & (0.048) \\[2pt]
R\&D intensity & 2.27$^{***}$ & 2.19$^{***}$ & 2.38$^{***}$ & 2.41$^{***}$ \\
 & (0.69) & (0.73) & (0.75) & (0.72) \\[2pt]
Leverage & $-$0.120 & $-$0.109 & $-$0.125 & $-$0.114 \\
 & (0.090) & (0.094) & (0.097) & (0.095) \\[2pt]
Slack & 0.187$^{*}$ & 0.211$^{**}$ & 0.179$^{*}$ & 0.193$^{*}$ \\
 & (0.097) & (0.101) & (0.105) & (0.102) \\[2pt]
Competition & $-$0.197$^{***}$ & $-$0.209$^{***}$ & $-$0.188$^{**}$ & $-$0.201$^{***}$ \\
 & (0.070) & (0.073) & (0.076) & (0.074) \\[2pt]
\midrule
Sample & Full & 2010 to 2019 & Excl. top size decile & Manufacturing-only \\
Firm fixed effects & Yes & Yes & Yes & Yes \\
Year fixed effects & Yes & Yes & Yes & Yes \\
Observations & 3,906 & 3,294 & 3,517 & 3,214 \\
R$^2$ & 0.095 & 0.098 & 0.093 & 0.101 \\
\bottomrule
\end{tabular}
\begin{minipage}{0.92\textwidth}\footnotesize\smallskip
\textit{Note.} Dependent variable: log of one plus new product introductions. Column 1 is the full model on the full sample. Column 4 keeps firms that draw most of their revenue from manufacturing. Standard errors, clustered by firm, in parentheses. \textasteriskcentered{} p$<$0.10, \textasteriskcentered{}\textasteriskcentered{} p$<$0.05, \textasteriskcentered{}\textasteriskcentered{}\textasteriskcentered{} p$<$0.01
\end{minipage}
\end{table}
\clearpage

\begin{table}[p]
\centering\small
\caption{Alliance Experience and Product Withdrawal: Cox Models (Hypothesis 5)}
\begin{tabular}{lcc}
\toprule
 & Model 1 & Model 2 \\
\midrule
Alliance experience &  & 0.18$^{***}$ \\
 &  & (0.06) \\[2pt]
Product novelty & 0.31$^{***}$ & 0.30$^{***}$ \\
 & (0.08) & (0.08) \\[2pt]
Firm size & $-$0.12$^{***}$ & $-$0.13$^{***}$ \\
 & (0.03) & (0.04) \\[2pt]
Firm age & $-$0.05 & $-$0.06 \\
 & (0.07) & (0.07) \\[2pt]
R\&D intensity & $-$1.12 & $-$1.21 \\
 & (1.24) & (1.25) \\[2pt]
\midrule
Industry strata & Yes & Yes \\
Products & 6,842 & 6,842 \\
Withdrawals & 1,107 & 1,107 \\
Log pseudolikelihood & $-$8,904.2 & $-$8,899.7 \\
\bottomrule
\end{tabular}
\begin{minipage}{0.92\textwidth}\footnotesize\smallskip
\textit{Note.} Cox proportional hazards models of the time from a product's introduction to its withdrawal. Coefficients are on the log-hazard scale: a positive coefficient means a higher withdrawal hazard, that is, a higher risk of withdrawal. Product novelty equals one when the product opens a product class that is new to the firm. Standard errors, clustered by firm, in parentheses. \textasteriskcentered{} p$<$0.10, \textasteriskcentered{}\textasteriskcentered{} p$<$0.05, \textasteriskcentered{}\textasteriskcentered{}\textasteriskcentered{} p$<$0.01
\end{minipage}
\end{table}
\clearpage

\begin{table}[p]
\centering\small
\caption{Alliance Experience and Process Innovation (Hypothesis 6)}
\begin{tabular}{lcc}
\toprule
 & Model 1 & Model 2 \\
\midrule
Alliance experience &  & 0.012 \\
 &  & (0.020) \\[2pt]
Firm size & 0.071$^{***}$ & 0.070$^{***}$ \\
 & (0.024) & (0.024) \\[2pt]
Firm age & 0.018 & 0.018 \\
 & (0.031) & (0.031) \\[2pt]
R\&D intensity & 1.62$^{***}$ & 1.60$^{***}$ \\
 & (0.48) & (0.48) \\[2pt]
Leverage & $-$0.044 & $-$0.044 \\
 & (0.061) & (0.061) \\[2pt]
Slack & 0.083 & 0.082 \\
 & (0.069) & (0.069) \\[2pt]
Competition & $-$0.072 & $-$0.072 \\
 & (0.048) & (0.048) \\[2pt]
\midrule
Firm fixed effects & Yes & Yes \\
Year fixed effects & Yes & Yes \\
Observations & 3,906 & 3,906 \\
R$^2$ & 0.052 & 0.052 \\
\bottomrule
\end{tabular}
\begin{minipage}{0.92\textwidth}\footnotesize\smallskip
\textit{Note.} Dependent variable: log of one plus new process filings. Standard errors, clustered by firm, in parentheses. \textasteriskcentered{} p$<$0.10, \textasteriskcentered{}\textasteriskcentered{} p$<$0.05, \textasteriskcentered{}\textasteriskcentered{}\textasteriskcentered{} p$<$0.01
\end{minipage}
\end{table}
\clearpage

\begin{table}[p]
\centering\small
\caption{Founder CEOs and the Effect of Alliance Experience (Hypothesis 8)}
\begin{tabular}{lcc}
\toprule
 & (1) & (2) \\
\midrule
Alliance experience & 0.139$^{***}$ & 0.151$^{***}$ \\
 & (0.041) & (0.042) \\[2pt]
Founder CEO & 0.052 & 0.118 \\
 & (0.061) & (0.074) \\[2pt]
Founder CEO $\times$ Alliance experience &  & $-$0.049$^{**}$ \\
 &  & (0.025) \\[2pt]
Firm size & 0.097$^{***}$ & 0.098$^{***}$ \\
 & (0.032) & (0.032) \\[2pt]
Firm age & $-$0.046 & $-$0.045 \\
 & (0.045) & (0.045) \\[2pt]
R\&D intensity & 2.28$^{***}$ & 2.27$^{***}$ \\
 & (0.69) & (0.69) \\[2pt]
Leverage & $-$0.121 & $-$0.120 \\
 & (0.090) & (0.090) \\[2pt]
Slack & 0.188$^{*}$ & 0.186$^{*}$ \\
 & (0.097) & (0.097) \\[2pt]
Competition & $-$0.198$^{***}$ & $-$0.197$^{***}$ \\
 & (0.070) & (0.070) \\[2pt]
\midrule
Firm fixed effects & Yes & Yes \\
Year fixed effects & Yes & Yes \\
Observations & 3,906 & 3,906 \\
R$^2$ & 0.095 & 0.097 \\
\bottomrule
\end{tabular}
\begin{minipage}{0.92\textwidth}\footnotesize\smallskip
\textit{Note.} Dependent variable: log of one plus new product introductions. Hypothesis 8 predicts negative moderation: a founder CEO weakens the positive association between alliance experience and new product introductions. Standard errors, clustered by firm, in parentheses. \textasteriskcentered{} p$<$0.10, \textasteriskcentered{}\textasteriskcentered{} p$<$0.05, \textasteriskcentered{}\textasteriskcentered{}\textasteriskcentered{} p$<$0.01
\end{minipage}
\end{table}
\clearpage

\begin{table}[p]
\centering\small
\caption{Alliance Experience in Family and Non-Family Firms (Hypothesis 9)}
\begin{tabular}{lcc}
\toprule
 & Family firms & Non-family firms \\
\midrule
Alliance experience & 0.18$^{**}$ & 0.09 \\
 & (0.07) & (0.08) \\[2pt]
Firm size & 0.13$^{***}$ & 0.10$^{***}$ \\
 & (0.04) & (0.03) \\[2pt]
Firm age & $-$0.03 & $-$0.06 \\
 & (0.07) & (0.06) \\[2pt]
R\&D intensity & 2.61$^{**}$ & 2.14$^{**}$ \\
 & (1.12) & (0.86) \\[2pt]
Leverage & $-$0.10 & $-$0.13 \\
 & (0.14) & (0.11) \\[2pt]
Slack & 0.16 & 0.19 \\
 & (0.15) & (0.14) \\[2pt]
Competition & $-$0.17 & $-$0.21$^{**}$ \\
 & (0.13) & (0.09) \\[2pt]
\midrule
Firm fixed effects & Yes & Yes \\
Year fixed effects & Yes & Yes \\
Firms & 136 & 276 \\
Observations & 1,288 & 2,618 \\
R$^2$ & 0.112 & 0.087 \\
\bottomrule
\end{tabular}
\begin{minipage}{0.92\textwidth}\footnotesize\smallskip
\textit{Note.} Dependent variable: log of one plus new product introductions. The two samples share no firms: a firm's family status does not change during the panel. Standard errors, clustered by firm, in parentheses. \textasteriskcentered{} p$<$0.10, \textasteriskcentered{}\textasteriskcentered{} p$<$0.05, \textasteriskcentered{}\textasteriskcentered{}\textasteriskcentered{} p$<$0.01
\end{minipage}
\end{table}
\clearpage

\begin{table}[p]
\centering\small
\caption{Board Size and New Product Introductions (Hypothesis 10)}
\begin{tabular}{lcccc}
\toprule
 & Model 1 & Model 2 & Model 3 & Model 4 \\
\midrule
Board size &  & 0.083$^{***}$ & 0.058$^{*}$ & 0.041 \\
 &  & (0.031) & (0.032) & (0.036) \\[2pt]
Alliance experience & 0.142$^{***}$ &  & 0.181$^{***}$ & 0.140$^{***}$ \\
 & (0.041) &  & (0.040) & (0.041) \\[2pt]
Firm size & 0.097$^{***}$ &  & 0.088$^{***}$ & 0.091$^{***}$ \\
 & (0.032) &  & (0.030) & (0.033) \\[2pt]
Firm age & $-$0.048 &  & $-$0.059 & $-$0.047 \\
 & (0.045) &  & (0.039) & (0.045) \\[2pt]
R\&D intensity & 2.29$^{***}$ &  & 2.38$^{***}$ & 2.27$^{***}$ \\
 & (0.69) &  & (0.66) & (0.69) \\[2pt]
Leverage & $-$0.121 &  & $-$0.126 & $-$0.119 \\
 & (0.090) &  & (0.085) & (0.090) \\[2pt]
Slack & 0.188$^{*}$ &  & 0.194$^{**}$ & 0.186$^{*}$ \\
 & (0.097) &  & (0.094) & (0.097) \\[2pt]
Competition & $-$0.198$^{***}$ &  & $-$0.203$^{***}$ & $-$0.196$^{***}$ \\
 & (0.070) &  & (0.066) & (0.070) \\[2pt]
\midrule
Firm fixed effects & Yes & No & No & Yes \\
Year fixed effects & Yes & Yes & Yes & Yes \\
Observations & 3,906 & 3,906 & 3,906 & 3,906 \\
R$^2$ & 0.094 & 0.011 & 0.103 & 0.095 \\
\bottomrule
\end{tabular}
\begin{minipage}{0.92\textwidth}\footnotesize\smallskip
\textit{Note.} Dependent variable: log of one plus new product introductions. Model 2 has no controls. Model 4 is the full model. Standard errors, clustered by firm, in parentheses. \textasteriskcentered{} p$<$0.10, \textasteriskcentered{}\textasteriskcentered{} p$<$0.05, \textasteriskcentered{}\textasteriskcentered{}\textasteriskcentered{} p$<$0.01
\end{minipage}
\end{table}
\clearpage

\begin{table}[p]
\centering\small
\caption{Absorptive Capacity as a Mediator (Hypothesis 11)}
\begin{tabular}{lccc}
\toprule
 & Model 1 & Model 2 & Model 3 \\
 & Intro. & Abs. cap. & Intro. \\
\midrule
Alliance experience & 0.30$^{***}$ & 0.42$^{***}$ & 0.08 \\
 & (0.08) & (0.09) & (0.07) \\[2pt]
Absorptive capacity &  &  & 0.51$^{***}$ \\
 &  &  & (0.12) \\[2pt]
Firm size & 0.13$^{***}$ & 0.33$^{***}$ & $-$0.05 \\
 & (0.04) & (0.06) & (0.05) \\[2pt]
Firm age & $-$0.06 & 0.09 & $-$0.08 \\
 & (0.05) & (0.07) & (0.07) \\[2pt]
Leverage & $-$0.15 & $-$0.21 & $-$0.04 \\
 & (0.10) & (0.14) & (0.10) \\[2pt]
Slack & 0.15 & 0.17 & 0.12 \\
 & (0.12) & (0.15) & (0.11) \\[2pt]
Competition & $-$0.24$^{***}$ & 0.04 & $-$0.24$^{***}$ \\
 & (0.07) & (0.10) & (0.08) \\[2pt]
\midrule
Industry fixed effects & Yes & Yes & Yes \\
Year fixed effects & Yes & Yes & Yes \\
Observations & 3,906 & 3,906 & 3,906 \\
R$^2$ & 0.214 & 0.388 & 0.262 \\
\bottomrule
\end{tabular}
\begin{minipage}{0.92\textwidth}\footnotesize\smallskip
\textit{Note.} Dependent variable in Models 1 and 3: log of one plus new product introductions; in Model 2: absorptive capacity (log R\&D staff). Pooled models with industry and year fixed effects; R\&D intensity is omitted because it is closely tied to R\&D staff. Standard errors, clustered by firm, in parentheses. \textasteriskcentered{} p$<$0.10, \textasteriskcentered{}\textasteriskcentered{} p$<$0.05, \textasteriskcentered{}\textasteriskcentered{}\textasteriskcentered{} p$<$0.01
\end{minipage}
\end{table}
\clearpage

\begin{table}[p]
\centering\small
\caption{Supplier Concentration and New Product Introductions (Hypothesis 12)}
\begin{tabular}{lcc}
\toprule
 & Model 1 & Model 2 \\
\midrule
Supplier concentration &  & 1.94$^{***}$ \\
 &  & (0.52) \\[2pt]
Supplier concentration squared &  & $-$1.25$^{***}$ \\
 &  & (0.38) \\[2pt]
Alliance experience & 0.142$^{***}$ & 0.139$^{***}$ \\
 & (0.041) & (0.041) \\[2pt]
Firm size & 0.097$^{***}$ & 0.094$^{***}$ \\
 & (0.032) & (0.032) \\[2pt]
Firm age & $-$0.048 & $-$0.049 \\
 & (0.045) & (0.045) \\[2pt]
R\&D intensity & 2.29$^{***}$ & 2.22$^{***}$ \\
 & (0.69) & (0.69) \\[2pt]
Leverage & $-$0.121 & $-$0.117 \\
 & (0.090) & (0.090) \\[2pt]
Slack & 0.188$^{*}$ & 0.183$^{*}$ \\
 & (0.097) & (0.097) \\[2pt]
Competition & $-$0.198$^{***}$ & $-$0.192$^{***}$ \\
 & (0.070) & (0.070) \\[2pt]
\midrule
Firm fixed effects & Yes & Yes \\
Year fixed effects & Yes & Yes \\
Observations & 3,906 & 3,906 \\
R$^2$ & 0.094 & 0.099 \\
\bottomrule
\end{tabular}
\begin{minipage}{0.92\textwidth}\footnotesize\smallskip
\textit{Note.} Dependent variable: log of one plus new product introductions. Standard errors, clustered by firm, in parentheses. \textasteriskcentered{} p$<$0.10, \textasteriskcentered{}\textasteriskcentered{} p$<$0.05, \textasteriskcentered{}\textasteriskcentered{}\textasteriskcentered{} p$<$0.01
\end{minipage}
\end{table}
\clearpage


\end{document}
